\section{Background}
\label{sec:background}

\paragraph{The DETR family.} DETR~\cite{carion2020detr} reframes detection as direct set prediction: a CNN backbone produces a feature map; a transformer encoder--decoder operates on a flattened token sequence; $N$ object queries cross-attend over encoder memory across $L$ decoder layers; a Hungarian matcher assigns predictions to ground-truth boxes for loss computation, removing the need for NMS. Deformable DETR~\cite{zhu2021deformable} replaces global attention with sparse, learned-offset attention over a multi-scale feature pyramid, dramatically speeding up convergence.

\paragraph{RT-DETR.} RT-DETR~\cite{lyu2024rtdetr} is a real-time variant of this family. Its encoder splits into two complementary blocks: AIFI, a transformer applied only to the smallest feature map ($S_5$, stride 32), and CCFF, a CNN-based fusion across pyramid levels. The flattened multi-level token sequence --- the \emph{encoder memory} --- is then attended over by a deformable decoder. Queries are initialized from the top-$N$ encoder predictions, in the spirit of DAB-DETR~\cite{liu2022dabdetr}.

\paragraph{Why small objects are hard.} Two structural issues compound. First, most published transformer detectors omit the stride-4 feature level (P2) for cost reasons; at stride 8, a $16$-pixel object covers only $2{\times}2$ tokens, and at stride 32 it is sub-token. Second, encoder memory is computed once before any decoder layer fires --- as the decoder discovers object hypotheses, that information cannot flow back to revise the features it is reading from. Refinement is strictly one-way.

We address the first issue by adding the P2 level (a standard FPN-style move~\cite{lin2017fpn}). Our novel contribution targets the second.
